\documentclass[11pt]{article}
\usepackage{lmodern}% scalable T1 font: without it pdfTeX falls back to Type3 bitmaps,
\usepackage{cmap}%     whose glyphs carry no Unicode mapping, so dashes and ligatures
%                       come out of the PDF text layer as raw encoding slots.
\usepackage[T1]{fontenc}
\usepackage[utf8]{inputenc}
\usepackage{mathptmx}
\usepackage[margin=1in]{geometry}
\usepackage[hidelinks]{hyperref}
\usepackage{enumitem}
\usepackage{microtype}
\usepackage{tabularx}
\usepackage[dvipsnames]{xcolor}
\definecolor{ReviewerGreen}{HTML}{008000}
\newcommand{\rvcomment}[1]{{\color{ReviewerGreen}#1}}
\newcommand{\concern}[1]{\textbf{\color{ReviewerGreen}#1}}
\setlength{\parindent}{0pt}
\setlength{\parskip}{6pt}
\setlist[itemize,1]{leftmargin=1.6em,labelsep=0.6em,itemsep=3pt,topsep=3pt,parsep=0pt,
                    label=\raisebox{-0.15ex}{\large\textbullet}}
\setlist[itemize,2]{leftmargin=1.4em,labelsep=0.5em,itemsep=2pt,topsep=3pt,parsep=0pt,
                    label=\textendash}
\pagestyle{plain}
\begin{document}

\vspace{12pt}
\concern{Editor, Concern \#1}

\concern{Editor comment:}

\rvcomment{``Please address the label-dependent aggregation rule by reporting primary results under a label-independent rule and providing sensitivity analyses under alternative tie policies. Also discuss the limitations of case construct validity and statistical independence assumptions.''}

\rvcomment{-Shiping Wang}

\textbf{Author response:}

\begin{itemize}

\item We agree. Preregistered decision~D4 is label-independent, with ties resolved toward \texttt{flag}. The submitted analysis instead applied strict majority to legitimate cases, an undisclosed deviation that we have corrected.

\item Under D4, false-positive rates are 2.3\%, 7.7\%, 37.0\%, 44.0\% and 85.7\%, giving an 83.3-percentage-point spread (95\% CI 71.8--93.7), versus 82.7 points as submitted. Across nine aggregation and single-call analyses, spreads range from 77.7 to 85.3 points without reversal of model ordering.

\item Synthetic labels were not independently adjudicated. To address case dependence, we added cluster-bootstrap intervals and a cluster-level permutation test (Monte Carlo $p \approx 0.0005$), alongside corrected case-level Cochran's $Q = 585.0$ ($p \approx 2.7\times10^{-125}$).

\end{itemize}

\textbf{Author action:}

\begin{itemize}

\item Section 2.3: Restored preregistered D4 and disclosed the deviation and its numerical consequences.

\item Section 2.4: Clarified independence assumptions, cluster-level inference and corrected Cochran's $Q$.

\item Section 4.2: Expanded construct-validity and statistical-independence limitations.

\item Appendix~A: Corrected the preregistered aggregation-rule description.

\item Table 1: Added aggregation and single-call sensitivity results, retaining the submitted-rule comparison.

\end{itemize}

\vspace{12pt}
\concern{Reviewer \#1, Concern \#1}

\concern{Reviewer comment:}

\rvcomment{The paper evaluates five foundation LLMs for AML screening, finding false-positive rates ranging from 0.3\% to 83.0\%.}

\rvcomment{Results are reported accurately with appropriate statistical measures.}

\rvcomment{Conclusions are well-founded with clear caveats about synthetic data limitations and hypothetical projections.}

\rvcomment{I recommend that the manuscript be accepted. The study is methodologically sound, results are reported accurately, and conclusions are appropriate.}

\textbf{Author response:}

\begin{itemize}

\item We thank Reviewer 1 for recommending acceptance.

\item The original 0.3\%--83.0\% false-positive range reflected the aggregation-rule deviation. Restoring D4 gives 2.3\%--85.7\%, and the central conclusion remains unchanged. The Cochran's $Q$ correction is detailed under Reviewer~2, Concern~\#6.

\end{itemize}

\vspace{12pt}
\concern{Reviewer \#2, Concern \#1}

\concern{Reviewer comment:}

\rvcomment{The paper asks whether the choice of commercial foundation model, rather than the screening task or prompt, materially changes the operating point of an LLM-based AML transaction screener. The authors run five zero-shot foundation models from two providers (three OpenAI, two Google) on an identical, preregistered battery of 600 synthetic transaction cases (300 suspicious across eight FATF-anchored typologies, 300 legitimate ``hard negative'' cases), crossed with two prompt formulations and two random seeds (12,000 calls total). They report a large spread in false-positive rate across models (0.3\%--83.0\%), a monotonic within-provider relationship between model tier and false-positive rate, high cross-model disagreement on legitimate cases (85.7\%), low chance-corrected agreement on misses, and an illustrative $\sim$197-fold difference in projected daily alert volume under a low-prevalence scenario. The paper frames these findings as a model-risk/governance issue rather than a model-capability issue, and supplies a full reproducibility package.}

\rvcomment{This is a carefully designed and quite transparent empirical study. The core idea --- treating foundation-model identity as the independent variable while holding the task and cases fixed, and reframing the result in model-risk-governance terms --- is a useful contribution to a literature that often focuses on identifying the best model. The statistical treatment is appropriate for the paired design, and the reproducibility package is appreciated.}

\rvcomment{I recommend minor revisions, focused mainly on (a) sharpening the originality claim (b) clarifying a few methodological choices that affect the headline numbers, and (c) some robustness check.}

\rvcomment{Please find the comments do be addressed in the uploaded file.}

\rvcomment{The central claim is credible but should be stated more precisely. The authors state they are ``the first study to control foundation-model identity as the experimental variable in AML transaction screening while fixing the task and transaction cases''. Within AML specifically this is plausible and, to my knowledge, not previously done in this controlled way. However, the general phenomenon being measured --- that different LLMs disagree substantially on identical inputs, and that prompt-induced variance is smaller than model-induced variance --- has precedent in the broader LLM-evaluation literature. I would ask the authors to sharpen the novelty claim to ``first domain-specific, governance-framed controlled measurement of this effect in AML screening'' rather than implying the phenomenon itself is new.}

\textbf{Author response:}

\begin{itemize}

\item We agree. Cross-model disagreement is established in broader LLM evaluation. Our contribution is its preregistered, governance-framed measurement in AML screening and translation into operational alert burden.

\item We also corrected the call count: 12,000 planned versus 12,120 archived calls, including 111 duplicate valid GPT-4o-mini calls and nine logged error calls. One GPT-4.1-mini case had only three valid calls.

\item Restricting voting to the four designed calls changes one GPT-4o-mini decision (FPR 85.7\% to 86.0\%, spread 83.3 to 83.7 points), without changing the conclusion.

\end{itemize}

\textbf{Author action:}

\begin{itemize}

\item Section 1: Narrowed the originality claim and cited broader LLM-evaluation work [26, 39].

\item Section 2.2: Documented realized call counts, retry handling and four-cell sensitivity.

\item Section 5: Scoped conclusions to the evaluated models and provider families.

\item Data availability: Updated the archived call count to 12,120.

\end{itemize}

\vspace{12pt}
\concern{Reviewer \#2, Concern \#2}

\concern{Reviewer comment:}

\rvcomment{The governance reframing is the paper's strongest and most original contribution. Treating model-tier substitution and provider-managed `-latest' aliases as first-class model-risk events, and proposing a standing ``operating-point probe'' as a monitoring control, is a useful and new operationalisation for this domain. This contribution stands on its own even if the underlying ``models disagree a lot'' finding were less novel, and I'd encourage the authors to foreground it even more.}

\textbf{Author response:}

\begin{itemize}

\item We agree that model governance deserves greater prominence. Model substitution and provider-managed alias redirection are now foregrounded as model-risk events, with a standing operating-point probe proposed for monitoring.

\end{itemize}

\textbf{Author action:}

\begin{itemize}

\item Section 1: Moved the governance contribution to first position.

\item Section 4.1: Added a dedicated operating-point-probe subsection covering triggers, monitored quantities and validation.

\item Abstract and Section 5: Foregrounded model substitution and alias-redirection risks.

\end{itemize}

\vspace{12pt}
\concern{Reviewer \#2, Concern \#3}

\concern{Reviewer comment:}

\rvcomment{A natural extension that bears on originality. The study covers only OpenAI and Google models, while the Generative AI declaration notes that Anthropic Claude was used for coding assistance. Given the paper's thesis is about provider dependence and model-swap risk, omitting other providers leaves the ``provider effect vs. tier effect'' distinction slightly underdetermined, a sentence explaining the provider selection (cost, access, or scope reasons) would help.}

\textbf{Author response:}

\begin{itemize}

\item The preregistered grid included xAI Grok-4.5. Decision D17 excluded Gemini Pro for thinking-mode incompatibility and cost, and D19 excluded Grok because credits were exhausted. The final five-model grid reflects documented attrition, limiting separation of provider and tier effects.

\item Anthropic Claude was used only for coding assistance, not evaluation.

\item We also corrected the Google metadata description: the archive records alias strings, not the concrete model versions served behind them.

\end{itemize}

\textbf{Author action:}

\begin{itemize}

\item Section 2.2: Documented the original three-provider grid and D17/D19 attrition.

\item Section 2.2: Clarified Claude's coding-only role and corrected Google alias metadata.

\item Section 4.2: Clarified provider--tier confounding and the need for third-provider replication.

\item Data availability: Specified that Google served-model fields contain aliases only.

\end{itemize}

\vspace{12pt}
\concern{Reviewer \#2, Concern \#4}

\concern{Reviewer comment:}

\rvcomment{Construct validity of the case battery. The false-positive/miss-rate estimates rest entirely on the correctness of the synthetic labels. The Methods section describes how cases were generated but does not state whether the ``legitimate hard negative'' and ``suspicious'' labels were reviewed or validated by an independent AML domain expert. I'd ask the authors to clarify whether any expert adjudication or inter-rater check was performed on a sample of the battery, and if not, to say so explicitly as a limitation.}

\textbf{Author response:}

\begin{itemize}

\item No independent AML expert adjudication or inter-rater assessment was performed. Labels were assigned by construction using FATF-anchored suspicious typologies and legitimate hard-negative templates.

\item Reported error rates therefore measure agreement with template-defined labels, not expert-adjudicated transaction risk. A model may flag a hard negative for reasons an analyst would also consider suspicious.

\item We did not substitute author adjudication for independent review. An expert-adjudicated subsample is identified as future work.

\end{itemize}

\textbf{Author action:}

\begin{itemize}

\item Section 2.1: Documented generator-defined labels and the absence of independent adjudication.

\item Section 4.2: Added construct-validity limitations and the need for expert-adjudicated evaluation.

\end{itemize}

\vspace{12pt}
\concern{Reviewer \#2, Concern \#5}

\concern{Reviewer comment:}

\rvcomment{The tie-breaking rule in the aggregation procedure (Appendix A, decision D4) deserves more scrutiny. The rule resolves 2--2 replicate splits ``with reference to the known case label'' (a tie counts as not a false positive for benign cases, and counted as caught for suspicious cases). This rule affects 132/1500 (8.8\%) of benign and 48/1500 (3.2\%) of suspicious model-case decisions, which is not negligible relative to some of the reported rates. I would ask for:}

\rvcomment{a breakdown of how many ties occurred per model, since if ties cluster in particular models, the label-referenced tie-break could differentially inflate or deflate that model's reported rate;}

\rvcomment{a sensitivity analysis reporting the headline spread (82.7 pp) and the cross-model Q-test under an alternative, label-agnostic tie rule, to show the qualitative conclusion is not an artifact of D4.}

\textbf{Author response:}

\begin{itemize}

\item Ties are unevenly distributed. Legitimate-case counts are 6, 20, 40, 58 and 8 across the five models, and suspicious-case counts are 19, 15, 10, 4 and 0. Restoring D4 increases false-positive rates most for the middle models, by 13.3 and 19.3 percentage points.

\item The spread changes from 82.7 to 83.3 points. Corrected Cochran's $Q$ is 585.0 (4 df, $p \approx 2.7\times10^{-125}$), versus 657.8 under the submitted rule. Gemini Flash versus GPT-4o also becomes significant after Bonferroni correction ($p = 0.004$).

\item Crucially, D4 was already preregistered as label-independent. The submitted legitimate-case rule was an undisclosed implementation deviation, now corrected and retained for sensitivity comparison.

\end{itemize}

\textbf{Author action:}

\begin{itemize}

\item Section 2.3: Restored D4 and documented the deviation and per-model tie effects.

\item Section 2.4: Corrected Cochran's $Q$ calculations.

\item Table 1: Reported sensitivity results across five aggregation rules and four single-call configurations.

\item Table 3: Updated paired McNemar comparisons under D4.

\item Table 5: Documented per-model tie counts and preregistered aggregation.

\end{itemize}

\vspace{12pt}
\concern{Reviewer \#2, Concern \#6}

\concern{Reviewer comment:}

\rvcomment{Statistical independence assumptions. McNemar's test and Cochran's Q assume the 300 legitimate cases are exchangeable/independent units. Since cases are synthetically generated from a shared typology structure, some correlation across cases seems plausible. The authors use cluster bootstrapping over the eight typology strata for the secondary uncertainty estimates but not, as far as I can tell, for the primary McNemar/Cochran tests. Given how strong the effect is ($P = 2\times10^{-27}$), this is unlikely to change the qualitative conclusion, but a brief comment on why case-level inference is appropriate for the primary test would be useful.}

\textbf{Author response:}

\begin{itemize}

\item We agree that cases sharing synthetic templates may be dependent. A cluster-level permutation test applying one model-label permutation per stratum gives Monte Carlo $p \approx 0.0005$ under either aggregation rule.

\item Cluster-bootstrap intervals for adjacent-tier differences are 29.3 points (95\% CI 12.6--51.2), 48.7 (25.8--72.5) and 41.7 (19.7--65.2). All exclude zero, although inference remains approximate with eight clusters.

\item The audit also identified an omitted factor of five in Cochran's $Q$ and chi-square tail-probability underflow. Corrected primary $Q = 585.0$ ($p \approx 2.7\times10^{-125}$), and submitted-rule $Q = 657.8$. Both errors were corrected without changing the qualitative conclusion or permutation ranking.

\end{itemize}

\textbf{Author action:}

\begin{itemize}

\item Section 2.4: Clarified dependence assumptions and cluster-level permutation inference.

\item Section 2.4: Documented the corrected $Q$ calculation and its independent verification.

\item Section 3.2: Added cluster-bootstrap intervals for adjacent-tier contrasts.

\item Reproducibility package: Corrected $Q$ and chi-square survival-function calculations, and regenerated and rehashed outputs.

\end{itemize}

\vspace{12pt}
\concern{Reviewer \#2, Concern \#7}

\concern{Reviewer comment:}

\rvcomment{``Model tier'' is treated ordinally but tested only descriptively. The paper observes a monotonic FPR decrease with tier within each provider and calls this a ``consistent pattern,''. I'd suggest adding a formal trend test to support the word ``monotonic,'' or otherwise soften ``monotonic'' to ``monotonically decreasing in the two families studied'' (which is closer to what is actually shown with only 2--3 tiers per provider).}

\textbf{Author response:}

\begin{itemize}

\item We added adjacent-tier paired McNemar tests rather than a Cochran--Armitage test, which would not respect the repeated-case design. All three contrasts remain significant after correction, with discordant counts of 2/90, 3/149 and 0/125.

\item The ordering holds within both evaluated provider families under the submitted rule, restored D4 and all four single-call configurations. Monotonicity claims are now restricted to these families.

\end{itemize}

\textbf{Author action:}

\begin{itemize}

\item Section 3.2: Added adjacent-tier paired tests and cluster-bootstrap intervals.

\item Section 3.2: Explained the paired-test choice and qualified monotonicity.

\item Abstract and Section 5: Restricted tier-ordering claims to the two evaluated provider families.

\end{itemize}

\vspace{12pt}
\concern{Reviewer \#2, Concern \#8}

\concern{Reviewer comment:}

\rvcomment{Operational projection is a single point, not a curve. The 197-fold alert-volume claim rests on one prevalence assumption (0.1\%) and one volume assumption (1,000,000/day). The paper is appropriately cautious that this is illustrative, but since alert volume and precision both scale nonlinearly with prevalence, a supplementary figure sweeping prevalence over a plausible range (e.g., 0.01\%--1\%) would let readers see whether the relative ordering of models (and the $\sim$200-fold spread) is stable across the plausible range, or whether it is sensitive to the specific 0.1\% choice.}

\textbf{Author response:}

\begin{itemize}

\item We swept suspicious-activity prevalence from 0.01\% to 1.0\%. Model ordering remains invariant, while the highest-to-lowest alert ratio under D4 is 36.6, 35.4 and 26.9 at 0.01\%, 0.1\% and 1.0\% respectively.

\item At 0.1\% prevalence, the submitted-rule ratio of approximately 197 falls to approximately 35 under D4, and the four single-call configurations yield 36- to 196-fold differences.

\item Projected precision ranges from approximately 0.4\%--27.6\% for Gemini Flash and 0.01\%--1.2\% for GPT-4o-mini across the sweep.

\end{itemize}

\textbf{Author action:}

\begin{itemize}

\item Section 3.3: Added prevalence-dependent alert-volume and precision comparisons.

\item Figure 4: Added the 0.01\%--1.0\% prevalence sweep and precision panel.

\item Abstract: Distinguished the 36- to 196-fold single-call range from the approximately 35-fold D4 result.

\end{itemize}

\vspace{12pt}
\concern{Reviewer \#2, Concern \#9}

\concern{Reviewer comment:}

\rvcomment{Zero-shot-only scope. This is a deliberate and well-justified design choice (isolating the base-model contribution before task-specific tuning), and the limitation is honestly stated. No change needed, but I'd suggest one added sentence in the Discussion speculating on whether calibration/threshold-tuning would be expected to compress the cross-model spread (since institutions could in principle recalibrate each model to a common operating point) or whether the paper's evidence (e.g., the persistence of ordering across two prompts) suggests the divergence is more structural than a simple decision-threshold shift. This matters for how alarming the governance implication should be considered.}

\textbf{Author response:}

\begin{itemize}

\item The evaluated interface returns binary decisions rather than calibrated scores, so threshold adjustment would require changing the elicitation procedure.

\item Model ordering persists across all four prompt--seed configurations, but this cannot establish whether score calibration would compress the spread. A scored evaluation with model-specific thresholds is needed to test that possibility.

\end{itemize}

\textbf{Author action:}

\begin{itemize}

\item Section 4: Clarified binary-output calibration limitations and configuration stability.

\item Section 4: Made the governance interpretation conditional on the evaluated elicitation procedure.

\end{itemize}

\vspace{12pt}
\concern{Reviewer \#2, Minor Point \#10}

\concern{Reviewer comment:}

\rvcomment{Please define ``model tier'' operationally in the Methods (currently it is defined via ``provider's contemporaneous capability/product and price positioning,'' which is reasonable but could usefully cite the specific list-price/capability source alongside the models.yaml registry reference).}

\textbf{Author response:}

\begin{itemize}

\item We clarified tier using within-provider price positioning. OpenAI list prices per million input/output tokens are \$2.50/\$10.00 for GPT-4o, \$0.40/\$1.60 for GPT-4.1-mini and \$0.15/\$0.60 for GPT-4o-mini.

\item Google's archived price fields are budgeting placeholders, not verified historical list prices, and the concrete versions behind the July 2026 aliases were not recorded.

\item Archived per-call expenditure corroborates within-provider ordering but also reflects token usage, and is not treated as a direct measure of capability or model size.

\end{itemize}

\textbf{Author action:}

\begin{itemize}

\item Section 2.4: Defined tier and cited verified OpenAI pricing.

\item Section 2.4: Distinguished Google budgeting placeholders from historical list prices.

\item Section 2.4: Added measured per-call costs and restricted comparisons to within-provider families.

\end{itemize}

\vspace{12pt}
\concern{Reviewer \#2, Minor Point \#11}

\concern{Reviewer comment:}

\rvcomment{The abstract states ``an 82.7 percentage-point spread'' while the Results also report a 95\% CI of 69.2--95.3 pp for that spread; consider including the CI in the abstract as well, since a $>$25-point CI width is itself informative.}

\textbf{Author response:}

\begin{itemize}

\item We agree. The revised abstract reports the D4 spread of 83.3 percentage points with its 95\% cluster-bootstrap confidence interval of 71.8--93.7 points, based on eight case-template strata.

\end{itemize}

\textbf{Author action:}

\begin{itemize}

\item Abstract: Added the corrected spread, confidence interval and cluster-bootstrap specification.

\end{itemize}

\vspace{12pt}
\concern{Reviewer \#2, Minor Point \#12}

\concern{Reviewer comment:}

\rvcomment{Consider explicitly stating in the Limitations that only two providers/five models were tested, so the paper's own headline ``foundation-model choice'' language should be read as scoped to the evaluated model set, not as a universal claim about all foundation models (the paper mostly does this correctly, but the Conclusion's phrasing is slightly more general than the evidence supports).}

\textbf{Author response:}

\begin{itemize}

\item We agree. Five models from two providers establish a substantial effect within the evaluated set, not a universal property of foundation models. Broader generalization requires replication.

\end{itemize}

\textbf{Author action:}

\begin{itemize}

\item Section 4.2: Limited generalizability to five models and two providers.

\item Section 5: Scoped the Conclusion and called for broader replication.

\end{itemize}

\vspace{12pt}
\concern{Reviewer \#2, Minor Point \#13}

\concern{Reviewer comment:}

\rvcomment{There are quite many typos such as numbers inserted across brackets in many paragraphs. Formatting also could be improved.}

\textbf{Author response:}

\begin{itemize}

\item We identified misplaced citations within em-dash asides and missing Unicode mappings for dashes and ligatures as sources of the formatting defects.

\item We corrected citation placement, standardized formatting and introduced a tested automated check for citation placement in the converted Word document.

\end{itemize}

\textbf{Author action:}

\begin{itemize}

\item Sections 1, 2.2 and 2.4: Corrected citation placement around em-dash asides.

\item Section 2.2: Replaced the design equation with prose.

\item Document-wide: Standardized numeric formatting, thousands separators and table alignment.

\item Word export: Added and tested automated citation-placement checks.

\item PDF export: Corrected Unicode mappings for dashes and ligatures.

\end{itemize}

\vspace{12pt}
\concern{Reviewer \#2, Recommendation}

\concern{Reviewer comment:}

\rvcomment{Minor revision. The core empirical contribution --- a controlled, paired, preregistered demonstration that foundation-model identity alone produces large, statistically robust shifts in AML screening operating points, reframed as a model-risk governance issue --- is original and well-executed. The requested revisions are addressable without new data collection, since the reproducibility package already contains everything needed to rerun the relevant analyses under alternative assumptions.}

\textbf{Author response:}

\begin{itemize}

\item We thank Reviewer 2 for recommending minor revision.

\end{itemize}

\vspace{12pt}
\concern{Reviewer \#3, Concern \#1}

\concern{Reviewer comment:}

\rvcomment{This manuscript investigates whether foundation-model choice materially changes the operating point of LLM-based anti-money-laundering transaction screening when the task and transaction cases are held constant. The topic is timely and the controlled cross-model design, paired statistical testing, reproducibility framework, and translation of false-positive rates into operational alert burden are valuable strengths. However, several methodological issues should be addressed before the conclusions can be fully supported, particularly the per-case aggregation procedure used to combine replicate model decisions.}

\rvcomment{Major comment --- label-dependent aggregation rule:}

\rvcomment{The preregistered aggregation rule resolves 2--2 replicate splits differently depending on the known ground-truth label: a tie on a legitimate case is treated as legitimate, whereas a tie on a suspicious case is treated as detected. Although this rule was preregistered and transparently disclosed, it uses information that would not be available in deployment and therefore may produce an optimistically biased evaluation operating point. This is especially important because ties affected a non-trivial number of model--case decisions. I recommend reporting the primary results under a label-independent aggregation rule and providing sensitivity analyses under alternative tie policies, such as treating ties as abstentions, always escalating ties, or analysing each prompt configuration separately. The authors should also clarify whether the headline false-positive and miss rates represent deployable model configurations or an evaluation-only aggregation procedure.}

\textbf{Author response:}

\begin{itemize}

\item We thank Reviewer 3 for identifying the central aggregation issue. The submitted label-dependent rule was an undisclosed deviation from preregistered D4, which specifies ties-to-flag without consulting ground truth. We restored D4 as primary and disclosed the deviation.

\item Sensitivity analyses confirm robustness: abstentions yield spreads of 84.9 points on model-specific retained cases and 81.5 on the common 189-case subset, ties-to-flag gives 83.3, unanimity 77.7, and the four single-call configurations 85.3, 85.0, 80.3 and 81.3 points. No model ordering reverses.

\item The four-call aggregate is an evaluation construct, with single-call results reported separately. The deviation understated false-positive rates without changing miss rates, and a label-independent strict majority reproduces the submitted false-positive rates but raises miss rates to 18.3\%, 8.3\%, 7.0\%, 2.0\% and 0.0\%.

\end{itemize}

\textbf{Author action:}

\begin{itemize}

\item Section 2.3: Restored D4 and disclosed the original deviation and its numerical consequences.

\item Section 2.3: Distinguished evaluation aggregation from operational deployment.

\item Table 1: Reported alternative aggregation and single-call sensitivity results.

\item Appendix~A: Corrected the D4 description and added the 189-case common-sample abstention analysis.

\item Table 1 and Appendix~A: Retained the submitted-rule comparison and clarified abstention denominators.

\end{itemize}

\end{document}
